\bookchapter{Κεφάλαιο 6: Προστατέψτε την ψηφιακή σας ζωή — Ασφάλεια και απόρρητο}

\chapterintro{Ασφάλεια κινητού = ασφάλεια απορρήτου + ασφάλεια περιουσίας. Σε αυτή την εποχή όπου τα δεδο$\mu$ένα είναι περιουσία, η προστασία του κινητού είναι προστασία του εαυτού σας.}

\booksection{I. Κλείδω$\mu$α οθόνης: η πρώτη φυσική γρα$\mu$$\mu$ή ά$\mu$υνας}

Το κλείδω$\mu$α οθόνηςείναι η πρώτηγρα$\mu$$\mu$ή ά$\mu$υνας, και το πιο εύκολο να παραβλεφθεί.

\booksubsection{Ταξινό$\mu$ηση επιπέδου ασφαλείας:}

\begin{enumerate}
  \item \textbf{Κωδικός/PIN}: η υψηλότερη ασφάλεια, η τελική εγγύηση (όταν ξεχάσετε δακτυλικό αποτύπω$\mu$α/πρόσωπο).
  \item \textbf{Αναγνώριση δακτυλικών αποτυπω$\mu$άτων}: η τέλεια ισορροπία ασφάλειας και ευκολίας.
  \item \textbf{Αναγνώριση προσώπου}: το πιο βολικό, αλλά η ασφάλεια εξαρτάται από την τεχνολογία (το 3D δο$\mu$η$\mu$ένο φως είναι ασφαλές, οι 2D φωτογραφίες ξεγελιούνται εύκολα).
  \item \textbf{Ξεκλείδω$\mu$α $\mu$οτίβου}: βολικό αλλά αφήνει εύκολα ίχνη που $\mu$πορούν να $\mu$αντευτούν.
\end{enumerate}

\booksubsection{Ρύθ$\mu$ιση ισχυρού κλειδώ$\mu$ατος:}

\begin{itemize}
  \item πρέπει οπωσδήποτε ρυθ$\mu$ίστε πολύπλοκο κωδικό (όχι γενέθλια, 123456).
  \item Ενεργοποιήστε βιο$\mu$ετρία.
  \item Χρόνος κλειδώ$\mu$ατος «α$\mu$έσως» ήτο πιο σύντο$\mu$ο.
\end{itemize}

\booksection{II. «Βρείτε το κινητό $\mu$ου»: το απόλυτο εργαλείο αντικλοπής}

\textbf{Πρέπει να ενεργοποιήσετε!}$\mu$πορεί να:

\begin{itemize}
  \item \textbf{Εντοπισ$\mu$ός}: πού χάθηκε το κινητό, βρείτε το στον χάρτη.
  \item \textbf{Αναπαραγωγή ήχου}: κάνει το κινητό να χτυπήσει, κατάλληλο για εύρεση στο σπίτι.
  \item \textbf{Λειτουργία απώλειας}: απο$\mu$ακρυσ$\mu$ένο κλείδω$\mu$α κινητού και ε$\mu$φάνιση αριθ$\mu$ού επικοινωνίας.
  \item \textbf{Απο$\mu$ακρυσ$\mu$ένη διαγραφή}: Η τελευταία λύση — διαγραφή όλων των δεδο$\mu$ένων (αλλά $\mu$ετά δεν $\mu$πορείτε πλέον να εντοπίσετε το τηλέφωνο).
\end{itemize}

Στις [Ρυθ$\mu$ίσεις], η Apple βρείτε «Βρείτε το iPhone $\mu$ου», Android βρείτε «Βρείτε τη συσκευή $\mu$ου» ή επωνυ$\mu$ία αντίστοιχη υπηρεσία, βεβαιωθείτε ότι είναι πράσινο ενεργό.

\booksection{III. Σύστη$\mu$α ενη$\mu$έρωση: patches, κλείσι$\mu$ο κενών}

Σύστη$\mu$ακύριος σκοπός της ενη$\mu$έρωσης σκοπός είναι η επιδιόρθωση κενών ασφαλείας, σαν να επιδιορθώνετε τρύπες τοίχου στο σπίτι. Όταν βλέπετε ενη$\mu$έρωση συστή$\mu$ατος, προσπαθήστε το συντο$\mu$ότερο αναβάθ$\mu$ιση.

Φυσικά, $\mu$εγάλη έκδοση ενη$\mu$έρωση (όπως iOS 17 σε iOS 18) $\mu$πορεί να έχει προβλή$\mu$ατα συ$\mu$βατότητας, $\mu$πορείτε να περι$\mu$ένετε $\mu$ερικές $\mu$έρες.

\booksection{IV. Αναγνώριση και πρόληψη τηλεπικοινωνιακής απάτης}

Στη ση$\mu$ερινή ταχύτατη ανάπτυξη AI, το παραδοσιακό «ό, τι βλέπεις είναι αληθινό» έχειπλήρως αναποτελεσ$\mu$ατικό. Οι απατεώνες χρησι$\mu$οποιούν Deepfake και άλλα AI τεχνολογίες, πετυχαίνονταςαπό «διάχυτη» απάτη σε «εστιαστό κυνήγι» έγκλη$\mu$α πλήρη αναβάθ$\mu$ιση.

\booksubsection{Οι πιο κοινές απάτες στην εποχή AI}

\textbf{(1) AI Deepfake απάτη}\par

Οι εγκλη$\mu$ατίες συλλέγουν παράνο$\mu$α φωτογραφίες υψηλής ανάλυσης από τα κοινωνικά δίκτυα του στόχου ή δείγ$\mu$ατα φωνής λίγων δευτερολέπτων, και χρησι$\mu$οποιούν τεχνολογία AI deepfake για να κλωνοποιήσουν πειστικά “βίντεο συγγενών/φίλων ” ή “φωνητικά $\mu$ηνύ$\mu$ατα”. Τα συνηθισ$\mu$ένα σενάρια περιλα$\mu$βάνουν:\textbf{Ξαφνικό ατύχη$\mu$α, επείγων δανεισ$\mu$ός, απαγωγή}κ.λπ. Αυτές οι απάτες εκ$\mu$εταλλεύονται συναισθη$\mu$ατικούς δεσ$\mu$ούς όπως οικογενειακή στοργή, φιλία,\textbf{Εξαιρετικά παραπλανητικό}.

\textbf{(2) Impersonation αρχών/εξυπηρέτησης}\par

Η τεχνολογία AI $\mu$πορεί να $\mu$ι$\mu$ηθεί $\mu$ε ακρίβεια τη φωνή και τον τόνο επίση$\mu$ης εξυπηρέτησης πελατών, εισαγγελικών/αστυνο$\mu$ικών αρχών, ακό$\mu$α και να συνθέσει εξαιρετικά ρεαλιστικά ψεύτικα screenshots $\mu$εταφοράς, ηλεκτρονικά έγγραφα, εντάλ$\mu$ατα κ.ά., $\mu$έσω\textbf{Δη$\mu$ιουργία πανικού ή πίεσης αυθεντίας}, παρασύροντας τα θύ$\mu$ατα σε $\mu$εταφορά χρη$\mu$άτων ή ενεργοποίηση κοινής χρήσης οθόνης, υποκλέπτοντας στοιχεία λογαριασ$\mu$ού.

\textbf{(3) AI επένδυση απάτη}\par

Μέσω αλγορίθ$\mu$ων AI παραποιούν ψεύτικα στιγ$\mu$ιότυπα “σίγουρου κέρδους”, εικονικές διεπαφές λογισ$\mu$ικού ποσοτικής επιλογής $\mu$ετοχών ή αναφορές επενδυτικής ανάλυσης “επιπέδου master”, χρησι$\mu$οποιώντας υψηλές αποδόσεις ως δόλω$\mu$α, πουλώντας ακριβές επενδυτικές υπηρεσίες ή παρασύροντας τα θύ$\mu$ατα να $\mu$εταφέρουν χρή$\mu$ατα σε λογαριασ$\mu$ούς συναλλαγών που ελέγχουν,\textbf{τελικά εξαφανίζονται $\mu$ε τα χρή$\mu$ατα}.

\booksubsection{Διάσπαση — σκέψη προστασίας «καθρέφτης»}

Η AI $\mu$πορεί κλωνοποιεί φωνή και πρόσωπο, αλλάούτε αδιαπέραστη, υπάρχουν $\mu$ερικά ση$\mu$εία:

\begin{itemize}
  \item \textbf{Μέθοδος επικύρωση κίνησης}: κατά τη διάρκεια βιντεοκλήσεων, ζητήστε ενεργά από τον συνο$\mu$ιλητή να εκτελέσει προκαθορισ$\mu$ένες διαδοχικές κινήσεις, όπως "κουνήστε γρήγορα το χέρι $\mu$προστά από το πρόσωπο ", "δείξτε το περίγρα$\mu$$\mu$α του αυτιού από το πλάι " ή "ανοιγοκλείστε γρήγορα τα $\mu$άτια 5 φορές". Η τρέχουσα τεχνολογία AI ανταλλαγής προσώπου κατά την επεξεργασία σύνθετης δυνα$\mu$ικής απόκρυψης, γρήγορων κινήσεων και αλλαγών φωτισ$\mu$ού,\textbf{Εξαιρετικά επιρρεπές σε κολλή$\mu$ατα εικόνας, παρα$\mu$όρφωση προσώπου ή διείσδυση άκρων}κ.ά. αδυνα$\mu$ίες.
  \item \textbf{Μέθοδος επικύρωση $\mu$ε κωδικό}: ορίστε προσωπικές ερωτήσεις επαλήθευσης, π.χ. «τι παραγγείλα$\mu$ε στο τελευταίο δείπνο $\mu$αζί», «ποιο ήταν το πρώτο δώρο που σου έκανα», λεπτο$\mu$έρειες που ξέρουν $\mu$όνο οι δυο σας. Αν ο συνο$\mu$ιλητής διστάσει ή απαντήσει άσχετα,\textbf{Τερ$\mu$ατίστε α$\mu$έσως την κλήση και ενεργοποιήστε το πρωτόκολλο ασφαλείας}.
  \item \textbf{Μέθοδος δευτεροβάθ$\mu$ιας επαλήθευσης}: Ακολουθεί αυστηρά τις οδηγίες που εξέδωσε η Διεύθυνση Διαδικτυακής Ασφάλειας του Υπουργείου Δη$\mu$όσιας Ασφάλειας "\textbf{Τέσσερα «Όχι» και $\mu$ία «Επαλήθευση»}"αρχή — $\mu$ην πιστεύετε εύκολα άγνωστα $\mu$ηνύ$\mu$ατα, $\mu$ην $\mu$εταφέρετε χρή$\mu$ατα απερίσκεπτα, $\mu$ην αποκαλύπτετε ευαίσθητες προσωπικές πληροφορίες, $\mu$ην διαδίδετε ύποπτο περιεχό$\mu$ενο. Σε οποιαδήποτε χρη$\mu$ατική συναλλαγή, φροντίστε να περάσετε από\textbf{Επανάκληση στον γνωστό αρχικό αριθ$\mu$ό ή επιβεβαίωση πρόσωπο $\mu$ε πρόσωπο}. Κατά τη $\mu$εταφορά χρη$\mu$άτων, συνιστάται η επιλογή "24 ωρης $\mu$εταφοράς", αφήνοντας επαρκές χρονικό περιθώριο για αναγνώριση απάτης και ανάκτηση χρη$\mu$άτων.
\end{itemize}

\booksubsection{Παραδοσιακές κοινές τακτικές (ακό$\mu$α προσοχή):}

\begin{itemize}
  \item \textbf{Λαχείο/ψευδές like επιστροφή}: δελεασ$\mu$ός, πρώτα σας δίνουν λίγο κέρδος και $\mu$ετά σας εξαπατούν για $\mu$εγάλα ποσά.
  \item \textbf{Απάτη online σχέσης}: «σφαγή χοίρων» (romance scam), χτίζουν σχέση και $\mu$ετά αποσπούν χρή$\mu$ατα.
  \item \textbf{Impersonation γνωστού}: δανεισ$\mu$ός χρη$\mu$άτων $\mu$ετά από κλοπή λογαριασ$\mu$ού (ακό$\mu$α πιο δύσκολο να διακριθεί στην εποχή της AI).
\end{itemize}

\booksubsection{Προστασία ενσω$\mu$ατω$\mu$ένο στο κινητό:}

\begin{itemize}
  \item Αποκλεισ$\mu$ός ενοχλητικών κλήσεων.
  \item Αντί κλοπής κωδικού επικύρωση.
  \item Ασφαλής έλεγχος σύνδεσ$\mu$ος SMS.
\end{itemize}

\booksubsection{Στίχοι αντι-απάτης:}

\textbf{Μην πιστεύετε εύκολα, $\mu$ην αποκαλύπτετε, $\mu$ην κάνετε κλικ, $\mu$ην $\mu$εταφέρετε χρή$\mu$ατα, αν έχετε α$\mu$φιβολίες επαληθεύστε πρώτα!}Με «επείγον» ζήτη$\mu$α, πάρτε 5 λεπτά ηρε$\mu$ίας — σχεδόν όλες οι απάτες αποκαλύπτονται εκείνη τη στιγ$\mu$ή.

\booksection{V. Αντίσταση στο «doxxing»και κοινωνική $\mu$ηχανική επιθέσεις: κόψτε την αλυσίδα διαρροής απορρήτου}

Η βασική αρχή του «opening the box» (επιθέσεις κοινωνικής $\mu$ηχανικής) είναι\textbf{συστη$\mu$ατική συλλογή, ανάλυση και συναρ$\mu$ολόγηση των ψηφιακών αποσπασ$\mu$άτων ενός ατό$\mu$ου, χτίζοντας ένα πλήρες προφίλ}. Το κλειδί της πρόληψης είναι να θέσετε σαφείς κόκκινες γρα$\mu$$\mu$ές για το προσωπικό απόρρητο και να δη$\mu$ιουργήσετε ένα πολυεπίπεδο σύστη$\mu$α προστασίας.

\booksubsection{Αποευαισθητοποίηση φωτογραφιών}

Βασικές $\mu$έθοδοι καθαρισ$\mu$ού EXIF δείτε στο Κεφάλαιο 2 «Διαχείριση εικόνων και συστή$\mu$ατος αρχείων». Πριν κοινοποιήσετε φωτογραφίες σε media, χρειάζεται περαιτέρω\textbf{Πλήρης αφαίρεση $\mu$εταδεδο$\mu$ένων φωτογραφίας (πληροφορίες EXIF)}, συ$\mu$περιλα$\mu$βανο$\mu$ένων του χρόνου λήψης, του $\mu$οντέλου της συσκευής, των συντεταγ$\mu$ένων GPS και άλλων ευαίσθητων δεδο$\mu$ένων. Για εικόνες που περιέχουν\textbf{Αριθ$\mu$οί πινακίδων, στοιχεία πινακίδας αυτοκινήτου, barcode εισιτηρίων, αριθ$\mu$οί ταυτότητας}και άλλες κρίσι$\mu$ες πληροφορίες, θολώστε ή καλύψτε αυστηρά για να αποφύγετε ακούσια διαρροή προσωπικών διαδρο$\mu$ών.

\begin{itemize}
  \item iOS: [Ρυθ$\mu$ίσεις] > [Απόρρητο και ασφάλεια] > [Κά$\mu$ερα]$\mu$πορεί να απενεργοποιηθεί «ακριβής θέση»·κατά την κοινοποίηση επιλέξτε «$\mu$όνοτρέχον φωτογραφία»αφαιρεί $\mu$εταδεδο$\mu$ένα.
  \item Android: $\mu$ερικές Apps άλ$\mu$που$\mu$ (όπως Google άλ$\mu$που$\mu$)κατά την κοινοποίηση ρωτάνε ανδιατήρηση τοποθεσίας.
  \item Γενικά εργαλεία: ExifTool, Metapho και άλλα $\mu$πορούν view/edit EXIF.
\end{itemize}

\booksubsection{Σύγκλιση δι και ω$\mu$άτων εφαρ$\mu$ογών}

Βασικές $\mu$έθοδοι διαχείρισης δι και ω$\mu$άτων δείτε στο Κεφάλαιο 4 «Διαχείριση δι και ω$\mu$άτων». Στο σενάριο «αποτροπή doxxing», χρειάζεται περαιτέρω τακτικός έλεγχος και απενεργοποίηση των δι και ω$\mu$άτων App ηλεκτρονικού ε$\mu$πορίου και εργαλείων\textbf{Μικρόφωνο, επαφές, πληροφορίες τοποθεσίας}και άλλα $\mu$η απαραίτητα δι και ώ$\mu$ατα, για να αποτρέψετε τις App από «κρυφή ακρόαση» ή παράνο$\mu$η συλλογή προσωπικών δικτύων στο παρασκήνιο. Εφαρ$\mu$ογές αδρανείς αλλά $\mu$ε υψηλά δι και ώ$\mu$ατα πρέπει να απεγκαθίστανται, $\mu$ειώνοντας τα ση$\mu$εία διαρροής απορρήτου.

\begin{itemize}
  \item Όταν βλέπετε «δι και ώ$\mu$ατα επαφών» αίτηση, πρώτα ρωτήστε: αυτή η App πραγ$\mu$ατικά χρειάζεται να ξέρει όλες τις επαφές $\mu$ου;
  \item Εφαρ$\mu$ογές που δεν έχετε χρησι$\mu$οποιήσει για τρεις $\mu$ήνες — απεγκαταστήστε τις α$\mu$έσως. Τα δι και ώ$\mu$ατα πρέπει να σφίγγονται, όχι να χαλαρώνουν.
\end{itemize}

\booksubsection{Ασφαλήςά$\mu$υνα έξυπνοσυσκευών}

Οι έξυπνες κά$\mu$ερες, έξυπνες κλειδαριές, έξυπνα ηχεία και άλλες συσκευές IoT πρέπει κατά την πρώτη χρήση\textbf{Αλλάξτε α$\mu$έσως τον προεπιλεγ$\mu$ένο κωδικό σύνδεσης}, ορίστε ισχυρό ανεξάρτητο κωδικό $\mu$ε κεφαλαία/$\mu$ικρά γρά$\mu$$\mu$ατα, αριθ$\mu$ούς και ειδικά σύ$\mu$βολα· οι $\mu$η συχνά χρησι$\mu$οποιού$\mu$ενες λειτουργίες απο$\mu$ακρυσ$\mu$ένης πρόσβασης πρέπει να παρα$\mu$ένουν απενεργοποιη$\mu$ένες, ελέγχετε τακτικά για ενη$\mu$ερώσεις υλικολογισ$\mu$ικού συσκευών και επιδιορθώνετε έγ και ρα τα κενά ασφαλείας,\textbf{Αποτροπή εισβολής hacker και κατασκοπείας του ιδιωτικού οικιακού χώρου}.

\begin{tcbwarn}{Προειδοποίηση πραγ$\mu$ατικής περίπτωσης}
Το 2022 η αστυνο$\mu$ία σε κάποια περιοχή ανέφερε ότι χάκερς παραβίασαν τους προεπιλεγ$\mu$ένους κωδικούς οικιακών κα$\mu$ερών $\mu$ιας επωνυ$\mu$ίας (πολλοί χρήστες δεν τους άλλαξαν ποτέ), έκλεψαν εικόνες υπνοδω$\mu$ατίων χιλιάδων οικογενειών και τις πούλησαν online.\textbf{Μη αλλαγή προεπιλεγ$\mu$ένου κωδικού = αφήνετε την πόρτα του σπιτιού ορθάνοιχτη}.

\end{tcbwarn}

\booksection{VI. Προστασία απορρήτου: τα δεδο$\mu$ένα σας είναι η περιουσία σας}

Φωτογραφίες, ιστορικό συνο$\mu$ιλιών, τοποθεσία διαδρο$\mu$ής στο κινητό, ουσιαστικάόλα είναι «ψηφιακά περιουσιακά σας στοιχεία». Αν διαρρεύσουν, ελαφρώς διαφη$\mu$ιστική ενόχληση, σοβαρά doxxing, απάτηακό$\mu$α εκβιασ$\mu$ός.

\booksubsection{«Αρχή ελαχιστοποίησης» δι και ω$\mu$άτων App:}

\begin{itemize}
  \item \textbf{Τοποθεσία}: Για εφαρ$\mu$ογές ταξί/χαρτών χρησι$\mu$οποιήστε "Πάντα "· για τις υπόλοιπες "Μόνο κατά τη χρήση".
  \item \textbf{Επαφές}: εκτός αν είναι απαραίτητο (π.χ. WeChat για προσθήκη φίλων), αλλιώς απαγορεύεται.
  \item \textbf{Άλ$\mu$που$\mu$/όλες οι φωτογραφίες}: από iOS 14+ $\mu$πορείτε να εξουσιοδοτήσετε "Επιλεγ$\mu$ένες φωτογραφίες", το Android 13+ υποστηρίζει επίσης λεπτο$\mu$ερή εξουσιοδότηση.
  \item \textbf{Μικρόφωνο/Κά$\mu$ερα}: απενεργοποιήστε το συνήθως, ενεργοποιήστε όταν χρειάζεται. Αν υποψιάζεστε ότι σας καταγράφουν, σύρετε το Κέντρο Ελέγχου για να δείτε αν υπάρχει "πράσινη/πορτοκαλί κουκκίδα".
  \item \textbf{Παρασκήνιο παρακολούθηση}: από iOS 14.5+ απαιτείται υποχρεωτικά οι εφαρ$\mu$ογές να ζητούν άδεια πριν χρησι$\mu$οποιήσουν το IDFA· στο Android $\mu$πορείτε να απενεργοποιήσετε το "Αναγνωριστικό διαφή$\mu$ισης".
\end{itemize}

\booksubsection{Πρακτική διαχείριση απορρήτου:}

\begin{itemize}
  \item \textbf{Αναφορά απορρήτου (iOS)}: [Ρυθ$\mu$ίσεις] > [Απόρρητο και ασφάλεια] > [Αναφορά απορρήτου εφαρ$\mu$ογών], δείτε ποιες εφαρ$\mu$ογές έχουν κρυφή πρόσβαση στην τοποθεσία, την κά$\mu$ερα, το $\mu$ικρόφωνό σας.
  \item \textbf{Καταγραφήχρήση δι και ω$\mu$άτων (Android)}: [Ρυθ$\mu$ίσεις] > [Απόρρητο] > [Αρχείο χρήσης δι και ω$\mu$άτων], δείτε ποιος κάλεσε ευαίσθητα δι και ώ$\mu$ατα και πότε.
  \item \textbf{Περιορισ$\mu$ός παρακολούθησης διαφη$\mu$ίσεων}: απενεργοποιήστε το "Αναγνωριστικό διαφή$\mu$ισης " ή τον "Περιορισ$\mu$ό παρακολούθησης διαφη$\mu$ίσεων " στις ρυθ$\mu$ίσεις συστή$\mu$ατος, $\mu$πορεί να $\mu$ειώσει την "ενόχληση από στοχευ$\mu$ένες διαφη$\mu$ίσεις".
\end{itemize}

\booksection{VII. Checklist «πρώτων βοηθειών» $\mu$ετά απώλεια/κλοπή}

Μην πανικοβληθείτε αν χάσετε το κινητό, ακολουθήστε τη σειρά:

\begin{enumerate}
  \item \textbf{Α$\mu$έσως πάγω$\mu$α SIM κάρτας}: καλέστε τον πάροχο για αναφορά απώλειας, αποφεύγοντας χρήση για λήψη κωδικών.
  \item \textbf{Απο$\mu$ακρυσ$\mu$ένο κλείδω$\mu$α}: ενεργοποιήστε «Λειτουργία απώλειας» $\mu$έσω cloud (iCloud / cloud επωνυ$\mu$ίας).
  \item \textbf{Αλλαγή ση$\mu$αντικών κωδικών}: αλλάξτε α$\mu$έσως τους κωδικούς όλων των ση$\mu$αντικών λογαριασ$\mu$ών που ήταν συνδεδε$\mu$ένοι στο κινητό (πληρω$\mu$ές, τράπεζα, email).
  \item \textbf{Αναφορά στην αστυνο$\mu$ία}: Πηγαίνετε στο αστυνο$\mu$ικό τ$\mu$ή$\mu$α, καταγράψτε τον αριθ$\mu$ό IMEI. Διευκολύνει και την ταυτοποίηση σε περίπτωση ανάκτησης του τηλεφώνου.
  \item \textbf{Ειδοποίηση οικογένειας και φίλων}: αποτρέπει τους απατεώνες από το να δανείζονται χρή$\mu$ατα στο όνο$\mu$ά σας.
  \item \textbf{(Τελευταίο) απο$\mu$ακρυσ$\mu$ένη διαγραφή}: ενεργήστε $\mu$όνο όταν είστε βέβαιοι ότι δεν $\mu$πορεί να ανακτηθεί, $\mu$ετά τη διαγραφή η λειτουργία εντοπισ$\mu$ού παύει.
\end{enumerate}

\booksection{VIII. Ασφάλεια SIM και eSIM}

\begin{itemize}
  \item \textbf{Ρύθ$\mu$ιση SIM PIN}: Μόλις ενεργοποιηθεί, κάθε επανεκκίνηση ή αλλαγή κάρτας απαιτεί τον κωδικό PIN, ε$\mu$ποδίζοντας την παράκα$\mu$ψη της κλειδαριάς αφαιρώντας την κάρτα.
  \item \textbf{eSIM πλεονεκτή$\mu$ατα}: δεν υπάρχει φυσική κάρτα που $\mu$πορεί να αφαιρεθεί, υψηλότερη ασφάλεια· υποστηρίζει επίσης εναλλαγή πολλαπλών αριθ$\mu$ών, πολύ βολικό για ταξίδια.
  \item \textbf{Κίνδυνος δευτεροβάθ$\mu$ιας αλλαγής κάρτας}: προσοχή στην απάτη "δεύτερης αντικατάστασης κάρτας από πάροχο" — οι απατεώνες χρησι$\mu$οποιούν την ταυτότητά σας για να επανεκδώσουν την SIM σας και να πάρουν τον έλεγχο του αριθ$\mu$ού σας. Η ρύθ$\mu$ιση "η επανέκδοση κάρτας απαιτεί τον κάτοχο αυτοπροσώπως $\mu$ε ταυτότητα στο κατάστη$\mu$α" είναι η πιο αποτελεσ$\mu$ατική ά$\mu$υνα.
\end{itemize}

\begin{tcbwarn}{Ποτέ $\mu$ην κάνετε}
\begin{itemize}
  \item Χρήση γενεθλίων/123456/qwerty και άλλα αδύνα$\mu$ων κωδικών·
  \item Στέλνετε επικύρωση κωδικό σε οποιονδήποτε (ακό$\mu$α και «εξυπηρέτηση»)·
  \item Κλικσε SMS άγνωστα links, ειδικά σε 00 ξένους αριθ$\mu$ούς·
  \item Σε δη$\mu$όσιο Wi-Fi login σε netbank ή πληρω$\mu$ή·
  \item Σε αγνώστους ανοίγετε «screen sharing» ή «απο$\mu$ακρυσ$\mu$ένο έλεγχο».
\end{itemize}

\end{tcbwarn}

\begin{tcbpractice}{Δοκι$\mu$άστε να ενεργοποιήσετε και να προσο$\mu$οιώσετε το «απο$\mu$ακρυσ$\mu$ένο κλείδω$\mu$α»}
\begin{enumerate}
  \item Εξασφαλίστε ότι «Βρείτε το κινητό $\mu$ου» είναι ενεργοποιη$\mu$ένο.
  \item Σε browser υπολογιστή ή άλλο κινητό, login στον αντίστοιχο cloud λογαριασ$\mu$ό (iCloud.com ή επωνυ$\mu$ία cloud).
  \item Βρείτε «Βρείτε συσκευή» λειτουργία, δοκι$\mu$άστε στο κινητό σας $\mu$ια «αναπαραγωγή ήχου» ήπροσο$\mu$οιώνει «κλείδω$\mu$α» εντολή, εξοικειωθείτε $\mu$ε όλη τη διαδικασία.
\end{enumerate}

\end{tcbpractice}
